\documentclass[10pt,a4paper]{amsart}
\usepackage[margin=19mm]{geometry}
\usepackage{amsmath,amssymb,mathtools}
\usepackage[scr=rsfs,cal=euler]{mathalfa}
\usepackage{xfrac}
\usepackage[unicode,hidelinks,bookmarks=false]{hyperref}
\newcommand{\ie}{i.e.\ }
\newcommand{\cf}{cf.\ }
\newcommand{\resp}{respectively\ }
\newcommand{\Subsection}{\subsection}
\providecommand{\nobreakdash}{\nobreak\hbox{-}}
\setlength{\emergencystretch}{2em}
\multlinegap0pt
\usepackage{siunitx}
\usepackage{siunitx}
\usepackage{siunitx}
\usepackage{siunitx}
\usepackage{mathtools}
\usepackage{amssymb}
\usepackage{algorithm}\usepackage{algpseudocode}
\usepackage{xcolor}
\usepackage{dsfont}
\usepackage[shortlabels]{enumitem}
\usepackage{siunitx}
\newtheorem{proposition}{Proposition}
\newcommand{\sigloss}{\mathcal{L}_{\text{sigmoid}}}
\title{Is Dimensionality a Barrier for Retrieval Models?}
\author{Kiril Bangachev, Guy Bresler, Jonathan Kogan, Yury Polyanskiy}
\date{Source: arXiv:2605.23556v1 (CC BY 4.0)}
\begin{document}
\maketitle

\noindent\emph{Source and licence.} Is Dimensionality a Barrier for Retrieval Models?. Version arXiv:2605.23556v1 (CC BY 4.0), distributed by arXiv under the Creative Commons Attribution 4.0 International licence, http://creativecommons.org/licenses/by/4.0/, as recorded in the arXiv licence metadata for this exact version and shown on the abstract page https://arxiv.org/abs/2605.23556v1. Full licence notice: \texttt{licenses/arxiv-2605.23556-CC-BY-4.0.txt}.\par\smallskip

\noindent\textit{Editorial scope.} Counted unit: the article's proposition prop:marginsigmoid (source0.tex lines 1563--1584). The article's complete original proof of it is delivered with it (source0.tex lines 1587--1715, one \texttt{proof} environment of 129 lines). No mathematics is written, repaired or re-derived: the statement and the proof are the article's own lines, in the article's own notation, and no retained line is substituted or re-flowed.  No further source-local context is needed: every cross-reference of every retained line resolves inside the two delivered spans.  Nothing was omitted from any retained span, so the package carries no omission record.  No retained line contains a citation, so the wrapper carries no bibliography and no bibliography entry.  No retained line contains a figure environment, an image-inclusion command or drawing code, so no image is delivered and the package declares no asset dependency.  Editorial formatting only, in addition: an AMS \texttt{amsart} preamble that restates the article's own declarations and macros used by the retained lines, namely the environments \texttt{proposition} on separate counters, together with the standard packages those lines need.  The retained environments print in this document as Proposition 1 in order of appearance, whereas the article prints the same items with the numbering of its own full document.  The complete native source of the article as distributed in the arXiv e-print for this exact version is bundled unchanged as \texttt{sources/201219/source0.tex}.
\begin{proposition}[Maximal-Margin Embeddings and Sigmoid Loss]
\label{prop:marginsigmoid} 
Consider the sigmoid loss with inverse-temperature $t>0$ and bias $b$ over embeddings $\{U_j\}_{j = 1}^N, \{V_i\}_{i = 1}^n:$ 
$$
\sigloss(\{U_j\}_{j = 1}^N, \{V_i\}_{i = 1}^n; t, b)\coloneqq \sum_{j\in [N], i \in [n]} \log\Big(1 + \exp\big((-1)^{A_{ji}}(t\langle U_j, V_i\rangle -b)\big)\Big).
$$
Then
\begin{enumerate}
    \item If $\sigloss(\{U_j\}_{j = 1}^N, \{V_i\}_{i = 1}^n; t, b) < \log 2,$ then there exists $m > 0$ such that $\{U_j\}_{j = 1}^N, \{V_i\}_{i = 1}^n$ is a margin-$m,$ relative-bias-$\frac{b}{t}$ embedding of $A$.
    \item If $\{U_j\}_{j=1}^N$ and $\{V_i\}_{i=1}^n$ form a margin-$m_*$, relative-bias-$\tau$ embedding of $A$ where
$$
m_* \coloneqq \min\left\{
\min_{A_{ji}=1}\left(\langle U_j,V_i\rangle-\tau\right),
\min_{A_{ji}=0}\bigl(\tau-\langle U_j,V_i\rangle\bigr)
\right\} > 0.
$$
Then
$$
\lim_{T\to\infty}\frac{1}{T}\log \sigloss(\{U_j\}_{j=1}^N,\{V_i\}_{i=1}^n;T,T\tau)=-m_*.
$$
\end{enumerate}
\end{proposition}

\begin{proof}[Proof of Proposition \ref{prop:marginsigmoid}]
% We begin with part (1). Assume that $\{U_j\}_{j=1}^N$ and $\{V_i\}_{i=1}^n$ form a margin-$m$, relative-bias-$\frac{b}{t}$ embedding of $A$. By definition, this means that $\langle U_j,V_i\rangle \ge \frac{b}{t}+m$ if $A_{ji} = 1$ and $\langle U_j,V_i\rangle \le \frac{b}{t}-m$ if $A_{ji}=0$. Now fix $i\in[N]$ and $j\in[n]$. If $A_{ji}=1$, then
% $$
% (-1)^{A_{ji}}\bigl(t\langle U_j,V_i\rangle-b\bigr)
% = -\bigl(t\langle U_j,V_i\rangle-b\bigr)
% \le -\Bigl(t\Bigl(\frac{b}{t}+m\Bigr)-b\Bigr)
% = -tm.
% $$
% If $A_{ji}=0$, then
% $$
% (-1)^{A_{ji}}\bigl(t\langle U_j,V_i\rangle-b\bigr)
% = t\langle U_j,V_i\rangle-b
% \le t\Bigl(\frac{b}{t}-m\Bigr)-b
% = -tm.
% $$
% Therefore every summand in $\sigloss(\{U_j\}_{j=1}^N,\{V_i\}_{i=1}^n;t,b)$ is at most $\log(1+e^{-tm})$. Summing over all $(i,j)\in[N]\times[n]$, we obtain
% $$
% \sigloss(\{U_j\}_{j=1}^N,\{V_i\}_{i=1}^n;t,b)\le Nn\log(1+e^{-tm}).
% $$
% Since $\log(1+e^{-tm})\to 0$ as $t\to\infty$, it follows that $Nn\log(1+e^{-tm})<\log(1+e^{-|b|})$ for all sufficiently large $t$. Hence $\sigloss(\{U_j\}_{j=1}^N,\{V_i\}_{i=1}^n;t,b)<\log(1+e^{-|b|})$ for all sufficiently large $t$.

We begin with part (1). 
Suppose that $\sigloss(\{U_j\}_{j=1}^N,\{V_i\}_{i=1}^n;t,b)<\log 2$. For each fixed pair $(j,i)$, the term $\log\Big(1 + \exp\big((-1)^{A_{ji}}(t\langle U_j, V_i\rangle -b)\big)\Big)$ is nonnegative. So, it follows that for every fixed pair $(j,i),$
$$
\log\Bigl(1+\exp\bigl((-1)^{A_{ji}}(t\langle U_j,V_i\rangle-b)\bigr)\Bigr)<\log 2.
$$
Since the function $x\mapsto \log(1+e^x)$ is strictly increasing, we obtain
$$
(-1)^{A_{ji}}(t\langle U_j,V_i\rangle-b)<0.
$$
If $A_{ji}=1$, then $-(t\langle U_j,V_i\rangle-b)<0$, so $t\langle U_j,V_i\rangle-b>0.$ If $A_{ji}=0$, then $t\langle U_j,V_i\rangle-b<0.$

Since there are finitely many pairs, set
$$
m\coloneqq
\min\left\{
\min_{A_{ji}=1}\left(\langle U_j,V_i\rangle-\frac {b}{t}\right),
\min_{A_{ji}=0}\left(\frac {b}{t}-\langle U_j,V_i\rangle\right)
\right\}>0.
$$
Then $\{U_j\}_{j=1}^N,\{V_i\}_{i=1}^n$ form a margin-$m$, relative-bias-$\frac {b}{t}$ embedding of $A$.



% Since there are only finitely many pairs $(i,j)$, we may set
% $$
% \eta\coloneqq \min\!\left(\min_{A_{ji}=1}\langle U_j,V_i\rangle,\min_{A_{ji}=0}\bigl(-\langle U_j,V_i\rangle\bigr)\right)>0.
% $$
% Now choose $t_0$ sufficiently large so that $\frac{|b|}{t}<\frac{\eta}{2}$ for all $t\ge t_0$. Then for all $t\ge t_0$, if $A_{ji}=1$, we have
% $$
% \langle U_j,V_i\rangle\ge \eta \ge \frac{b}{t}+\frac{\eta}{2},
% $$
% and if $A_{ji}=0$, we have
% $$
% \langle U_j,V_i\rangle\le -\eta \le \frac{b}{t}-\frac{\eta}{2}.
% $$
% Therefore, for all sufficiently large $t$, $\{U_j\}_{j=1}^N$ and $\{V_i\}_{i=1}^n$ form a margin-$\frac{\eta}{2}$, relative-bias-$\frac{b}{t}$ embedding of $A$.

We now prove part 2. For each pair $(j,i) \in [N] \times [n]$, define
$$
\gamma_{ji}\coloneqq
\begin{cases}
\langle U_j,V_i\rangle-\tau, & A_{ji}=1,\\[4pt]
\tau-\langle U_j,V_i\rangle, & A_{ji}=0.
\end{cases}
$$
Then $\gamma_{ji}>0$ for all $(j,i)$, and by definition $m_*=\min_{j,i}\gamma_{ji}$.

If $A_{ji}=1$, then
$$
(-1)^{A_{ji}}\bigl(T\langle U_j,V_i\rangle-T\tau\bigr)
=
-\bigl(T\langle U_j,V_i\rangle-T\tau\bigr)
=
-T(\langle U_j,V_i\rangle-\tau)
=
-T\gamma_{ji}.
$$
If $A_{ji}=0$, then
$$
(-1)^{A_{ji}}\bigl(T\langle U_j,V_i\rangle-T\tau\bigr)
=
T\langle U_j,V_i\rangle-T\tau
=
-T(\tau-\langle U_j,V_i\rangle)
=
-T\gamma_{ji}.
$$
Therefore,
$$
\sigloss(\{U_j\}_{j=1}^N,\{V_i\}_{i=1}^n;T,T\tau)
=
\sum_{j\in[N],\,i\in[n]}\log\!\bigl(1+e^{-T\gamma_{ji}}\bigr).
$$

Since $\gamma_{ji}\ge m_*$ for all $(i,j)$, we obtain the upper bound
$$\sigloss(\{U_j\}_{j = 1}^N,\{V_i\}_{i = 1}^n;T,T\tau)
\leq
Nn\log(1+e^{-Tm_*})
\leq
Nne^{-Tm_*},
$$
where we used $\log(1+x)\le x$ for all $x\ge0$. For the lower bound, choose $(j_0,i_0)$ such that $\gamma_{j_0i_0}=m_*$. Then $\sigloss(\{U_j\}_{j=1}^N,\{V_i\}_{i = 1}^n;T,T\tau)
\ge
\log(1+e^{-Tm_*})$. Also, for $x\in[0,1]$, $\log(1+x)\ge \frac{x}{2}$. Hence for all sufficiently large $T$ (so that $e^{-Tm_*}\le1$), $\log(1+e^{-Tm_*})\ge \frac12 e^{-Tm_*}$. Thus for all sufficiently large $T$,
$$
\frac{1}{2} e^{-Tm_*}
\leq
\sigloss(\{U_j\}_{j = 1}^N,\{V_i\}_{i = 1}^n;T,T\tau)
\le
Nne^{-Tm_*}.
$$

Taking logarithms gives
$$
-Tm_*-\log 2
\le
\log \sigloss(\{U_j\}_{j = 1}^N,\{V_i\}_{i = 1}^n;T,T\tau)
\le
-Tm_*+\log(Nn).
$$
Dividing by $T$ and letting $T\to\infty$, we conclude that
$$
\lim_{T\to\infty}
\frac{1}{T}\log \sigloss(\{U_j\}_{j = 1}^N,\{V_i\}_{i = 1}^n;T,T\tau)
=
-m_*\qedhere
$$
\end{proof}

\end{document}
